\documentclass[11pt]{amsart}
\usepackage{amsmath,amssymb,amsthm}
\usepackage[margin=1in]{geometry}
\newtheorem{theorem}{Theorem}[section]
\newtheorem{lemma}[theorem]{Lemma}
\newtheorem{proposition}[theorem]{Proposition}
\theoremstyle{remark}\newtheorem{remark}[theorem]{Remark}
\newcommand{\E}{\mathbb{E}}
\title[The smallest NTK eigenvalue in the ordered phase]{The smallest eigenvalue of the neural tangent kernel of a deep tanh network in the ordered phase}
\author{Yitzchak Shmalo}
\begin{document}
\begin{abstract}
For a fully connected tanh network of depth $L$ in the infinite-width limit, with weight variance $\sigma_w^2$ and bias
variance $\sigma_b^2>0$ in the ordered phase $\chi_1<1$, we determine the smallest eigenvalue of the neural tangent kernel
Gram matrix on $n$ distinct inputs of the sphere. It equals
$L\chi_1^L\,(q^*+B\,T^*/\chi_1)\,\lambda_{\min}^{\perp}(-A)\,(1+O(1/L))$, where $q^*$ is the fixed-point preactivation
variance, $T^*=q^*/(1-\chi_1)$, $B=\chi'(1)$, and $A$ is the matrix of limiting rescaled correlation gaps. That the
non-constant eigenvalues are of order $L\chi_1^L$ is due to Xiao, Pennington and Schoenholz; the contribution here is the
constant, its positivity, and the extension to inputs of unequal norm. The factor $L$ comes from a double root in the
linear recursion for the kernel gap. We prove that $\lambda_{\min}^{\perp}(-A)>0$ for any
distinct inputs by showing that every Taylor coefficient of the limiting gap profile is strictly positive, so the
rate cannot degenerate further.
\end{abstract}
\maketitle

\section{Setting and main result}
Let $\phi=\tanh$. A fully connected network of depth $L$ with weights of variance $\sigma_w^2/n_{\ell}$ and biases of
variance $\sigma_b^2$ has, in the limit of infinite width, preactivations that are centred Gaussian processes with
covariance given by the recursion
\[
\Sigma_{\ell+1}(x,y)=\sigma_w^2\,\E\big[\phi(u)\phi(v)\big]+\sigma_b^2,\qquad
(u,v)\sim N\!\Big(0,\begin{pmatrix}\Sigma_\ell(x,x)&\Sigma_\ell(x,y)\\ \Sigma_\ell(x,y)&\Sigma_\ell(y,y)\end{pmatrix}\Big),
\]
and the neural tangent kernel satisfies
\[
\Theta_{\ell+1}(x,y)=\Sigma_{\ell+1}(x,y)+\dot\Sigma_{\ell+1}(x,y)\,\Theta_\ell(x,y),\qquad
\dot\Sigma_{\ell+1}(x,y)=\sigma_w^2\,\E\big[\phi'(u)\phi'(v)\big].
\]
We assume $\sigma_b^2>0$. The diagonal $q_\ell=\Sigma_\ell(x,x)$ obeys $q_{\ell+1}=\sigma_w^2\E[\phi(\sqrt{q_\ell}\,z)^2]+\sigma_b^2$
and converges to a unique fixed point $q^*>0$. At $q^*$ write $\chi(c)=\sigma_w^2\E[\phi'(u)\phi'(v)]$ for $(u,v)$ of
variance $q^*$ and correlation $c$, and
\[
\chi_1=\chi(1)=\sigma_w^2\,\E\big[\phi'(\sqrt{q^*}z)^2\big],\qquad B=\chi'(1),\qquad T^*=\frac{q^*}{1-\chi_1}.
\]
The ordered phase is $\chi_1<1$. In it every correlation $c_\ell(x,y)=\Sigma_\ell(x,y)/\sqrt{q_\ell(x)q_\ell(y)}$ tends
to $1$, and
\[
A_{ij}=\lim_{\ell\to\infty}\chi_1^{-\ell}\big(1-c_\ell(x_i,x_j)\big)
\]
exists (Section 2). Let $\mathbf 1^\perp=\{v\in\mathbb R^n:\sum_i v_i=0\}$ and
$\lambda_{\min}^\perp(M)=\min\{v^\top Mv: v\in\mathbf 1^\perp,\ |v|=1\}$.

\begin{theorem}\label{thm:main}
Let $x_1,\dots,x_n$ be distinct points with $\Sigma_0(x_i,x_i)=q^*$, and let $\Theta_L$ be the $n\times n$ Gram matrix
of the neural tangent kernel at depth $L$. Then, as $L\to\infty$,
\[
\lambda_{\min}(\Theta_L)=\Big(q^*+\frac{B\,T^*}{\chi_1}\Big)\,\lambda_{\min}^\perp(-A)\;L\,\chi_1^{L}\,\big(1+O(1/L)\big),
\]
and $\lambda_{\min}^\perp(-A)>0$. For distinct inputs of arbitrary positive norm the same asymptotic form holds, with
$A$ defined through the correlations of the network started at those inputs, whenever $\lambda_{\min}^\perp(-A)>0$.
\end{theorem}

\begin{remark}
The rate is $L\chi_1^L$ rather than $\chi_1^L$ because the gap $\Theta_\ell(x,x)-\Theta_\ell(x,y)$ satisfies a linear
recursion whose homogeneous rate and forcing rate are both $\chi_1$; the double root produces the factor $L$.
\end{remark}

For inputs normalized to $q^*$, Xiao, Pennington and Schoenholz \cite{xps} showed that the deviation of $\Theta_L$ from its limit,
rescaled by $L^{-1}\chi_1^{-L}$, converges and that the eigenvalues of $\Theta_L$ other than the largest are
$O(L\chi_1^L)$. Theorem~\ref{thm:main} identifies the limit as $\kappa(-A)$ with $\kappa=q^*+BT^*/\chi_1$
(Section~\ref{sec:ninput}), proves that $-A$ is positive definite on $\mathbf 1^\perp$, so that the order is attained by
the smallest eigenvalue and not only bounded above (Section~\ref{sec:positive}), and extends the asymptotics to inputs of
unequal norm (Section~\ref{sec:offnorm}).

\section{The correlation map and the double root}\label{sec:double}
Fix two inputs $x,y$ with $q_0(x)=q_0(y)=q^*$, so that $q_\ell=q^*$ at every layer, and write $c_\ell=c_\ell(x,y)$ and
$D_\ell=1-c_\ell$. The correlations obey $c_{\ell+1}=C(c_\ell)$ with
\[
C(c)=\frac{1}{q^*}\Big(\sigma_w^2\,\E\big[\phi(u)\phi(v)\big]+\sigma_b^2\Big),
\]
$(u,v)$ of variance $q^*$ and correlation $c$. Differentiating under the Gaussian integral gives $C'(c)=\chi(c)$, and
expanding $\phi'$ in Hermite polynomials shows that $\chi(c)=\sum_{k\ge0}b_kc^k$ with $b_k\ge0$. Hence $C$ is convex on
$[-1,1]$, $C(1)=1$ and $C'(1)=\chi_1$.

\begin{lemma}\label{lem:cmap}
In the ordered phase, $D_\ell\to0$ and the limit $A=\lim_{\ell\to\infty}\chi_1^{-\ell}D_\ell$ exists, with $A>0$ when
$c_0<1$ and
\[
D_\ell=A\chi_1^{\ell}\big(1+O(\chi_1^{\ell})\big).
\]
\end{lemma}

\begin{proof}
Convexity and $C'(1)=\chi_1<1$ give $C(c)>c$ for $c<1$, so $c_\ell$ increases to the only fixed point in $[-1,1]$, which is
$1$. Taylor's theorem at $c=1$ gives
\[
D_{\ell+1}=\chi_1D_\ell-\tfrac{B}{2}D_\ell^2+O(D_\ell^3),\qquad B=\chi'(1)\ge0,
\]
so $D_{\ell+1}/D_\ell=\chi_1(1+O(D_\ell))$. Since $D_\ell$ decays geometrically, $\sum_\ell D_\ell<\infty$ and the product
$\chi_1^{-\ell}D_\ell=D_0\prod_{k<\ell}\chi_1^{-1}D_{k+1}/D_k$ converges to a positive limit, the tail of the product being
$1+O(\sum_{k\ge\ell}D_k)=1+O(\chi_1^\ell)$.
\end{proof}

The same argument applied to each pair defines the matrix $A_{ij}$ of Theorem~\ref{thm:main}, with $A_{ii}=0$.

Turn to the kernel. Take $\Theta_0=\Sigma_0$ and put $T_\ell=\Theta_\ell(x,x)$. On the diagonal $T_{\ell+1}=q^*+\chi_1T_\ell$,
so $T_\ell=T^*+O(\chi_1^\ell)$. Off the diagonal $\Sigma_{\ell+1}(x,y)=q^*c_{\ell+1}$ and $\dot\Sigma_{\ell+1}(x,y)=\chi(c_\ell)$.
The gap $E_\ell=T_\ell-\Theta_\ell(x,y)$ therefore satisfies
\begin{equation}\label{eq:gap}
E_{\ell+1}=\chi_1E_\ell+q^*D_{\ell+1}+\big(\chi_1-\chi(c_\ell)\big)\,\Theta_\ell(x,y).
\end{equation}

\begin{proposition}\label{prop:two}
With $A$ as in Lemma~\ref{lem:cmap},
\[
E_L=A\Big(q^*+\frac{B\,T^*}{\chi_1}\Big)L\chi_1^L\big(1+O(1/L)\big).
\]
\end{proposition}

\begin{proof}
By Lemma~\ref{lem:cmap}, $q^*D_{\ell+1}=q^*A\chi_1^{\ell+1}(1+O(\chi_1^\ell))$ and $\chi_1-\chi(c_\ell)=BD_\ell+O(D_\ell^2)$.
A first pass through \eqref{eq:gap} with the crude bound $|\Theta_\ell(x,y)|\le T_\ell$ gives $E_\ell=O(\ell\chi_1^\ell)$,
hence $\Theta_\ell(x,y)=T_\ell-E_\ell=T^*+O(\ell\chi_1^\ell)$. The forcing in \eqref{eq:gap} is then
\[
f_\ell=A\big(q^*\chi_1+BT^*\big)\chi_1^\ell\big(1+O(\ell\chi_1^\ell)\big),
\]
and solving the recursion,
\[
E_L=\chi_1^LE_0+\sum_{\ell=0}^{L-1}\chi_1^{L-1-\ell}f_\ell
=A\big(q^*\chi_1+BT^*\big)\chi_1^{L-1}\Big(L+O\Big(\sum_{\ell\ge0}\ell\chi_1^\ell\Big)\Big)+O(\chi_1^L).
\]
The sum is finite, which gives the claim.
\end{proof}

For two inputs the Gram matrix has equal diagonal entries, so its smaller eigenvalue is exactly $E_L$, and
Proposition~\ref{prop:two} is Theorem~\ref{thm:main} for $n=2$, where $\lambda_{\min}^\perp(-A)=A_{12}$.

\section{The $n$-input Gram matrix}\label{sec:ninput}
Let $x_1,\dots,x_n$ be distinct with $\Sigma_0(x_i,x_i)=q^*$. Distinct inputs of equal norm have $c_0(x_i,x_j)<1$, so
Lemma~\ref{lem:cmap} gives $A_{ij}>0$ for $i\ne j$, and $A_{ii}=0$. Put $\kappa=q^*+BT^*/\chi_1$ and let $J$ be the
all-ones matrix. Every diagonal entry of $\Theta_L$ equals $T_L$ and, by Proposition~\ref{prop:two} applied to each of the
finitely many pairs,
\begin{equation}\label{eq:gram}
\Theta_L=T_LJ+\kappa L\chi_1^L(-A)+R_L,\qquad \|R_L\|=O(\chi_1^L).
\end{equation}

\begin{proposition}\label{prop:ninput}
$\lambda_{\min}(\Theta_L)=\kappa\,\lambda_{\min}^\perp(-A)\,L\chi_1^L+O(\chi_1^L)$.
\end{proposition}

\begin{proof}
For the upper bound take a unit $u\in\mathbf 1^\perp$ with $u^\top(-A)u=\lambda_{\min}^\perp(-A)$. Since $Ju=0$,
\eqref{eq:gram} gives $u^\top\Theta_Lu=\kappa\lambda_{\min}^\perp(-A)L\chi_1^L+O(\chi_1^L)$.

For the lower bound write a unit vector as $w=\alpha n^{-1/2}\mathbf 1+\beta u$ with $u\in\mathbf 1^\perp$ a unit vector and
$\alpha^2+\beta^2=1$. Then $w^\top Jw=n\alpha^2$, and with $\varepsilon_L=\kappa L\chi_1^L$,
\[
w^\top\Theta_Lw\ \ge\ nT_L\alpha^2-2\|A\|\,\varepsilon_L|\alpha\beta|+\beta^2\varepsilon_L\lambda_{\min}^\perp(-A)-O(\chi_1^L).
\]
By $2ab\le a^2+b^2$, $2\|A\|\varepsilon_L|\alpha\beta|\le\tfrac12nT_L\alpha^2+2\|A\|^2\varepsilon_L^2\beta^2/(nT_L)$, and
the last term is $O(L^2\chi_1^{2L})=O(\chi_1^L)$. So
\[
w^\top\Theta_Lw\ \ge\ \tfrac12nT_L\alpha^2+\beta^2\varepsilon_L\lambda_{\min}^\perp(-A)-O(\chi_1^L)
\ \ge\ \min\big(\tfrac12nT_L,\ \varepsilon_L\lambda_{\min}^\perp(-A)\big)-O(\chi_1^L),
\]
using $\lambda_{\min}^\perp(-A)\ge0$ (shown below). For large $L$ the minimum is $\varepsilon_L\lambda_{\min}^\perp(-A)$.
\end{proof}

The constant direction carries the eigenvalue $nT^*+O(L\chi_1^L)$, which stays of order one; every other eigenvalue
decays like $L\chi_1^L$, at rates set by the spectrum of $-A$ on $\mathbf 1^\perp$. Once
$\lambda_{\min}^\perp(-A)>0$ is known (Section~\ref{sec:positive}), Proposition~\ref{prop:ninput} is
Theorem~\ref{thm:main} for inputs on the fixed-point sphere.

That $-A$ is at least nonnegative on $\mathbf 1^\perp$ is elementary. For each $\ell$ the correlations $c_\ell(x_i,x_j)$
form a positive semidefinite matrix $K_\ell$ with unit diagonal, so $1-c_\ell(x_i,x_j)=\tfrac12|z_i-z_j|^2$ for some
unit vectors $z_i$, and for $v\in\mathbf 1^\perp$,
\[
v^\top\big(-(J-K_\ell)\big)v=v^\top K_\ell v\ge0 .
\]
Dividing by $\chi_1^\ell$ and letting $\ell\to\infty$ gives $v^\top(-A)v\ge0$. Strictness needs more, because a limit of
positive definite matrices can be singular.

\section{Strict positivity of $\lambda_{\min}^{\perp}(-A)$}\label{sec:positive}
Write the correlation map as a power series. Expanding $\phi(\sqrt{q^*}z)=\sum_kb_kH_k(z)$ in Hermite polynomials,
Mehler's formula gives
\[
C(c)=\sum_{k\ge0}a_kc^k,\qquad a_0=\frac{\sigma_b^2}{q^*},\qquad a_k=\frac{\sigma_w^2\,k!\,b_k^2}{q^*}\ (k\ge1),
\]
using $b_0=0$, since $\tanh$ is odd. Thus $a_k\ge0$, $\sum_ka_k=C(1)=1$, $a_0>0$, and $a_1>0$ because $\tanh$ is
increasing. Infinitely many $a_k$ are positive, since $\tanh$ is not a polynomial. In particular $|C(c)|\le C(|c|)\le1$
on $[-1,1]$. Put
\[
\Phi_\ell(c)=\chi_1^{-\ell}\big(C^{\circ\ell}(c)-1\big),\qquad c\in[-1,1],
\]
so that $\Phi_\ell(c_0)=-\chi_1^{-\ell}D_\ell$ and, by Lemma~\ref{lem:cmap}, $\Phi_\ell\to\Phi$ pointwise with
$\Phi(c_0(x_i,x_j))=-A_{ij}$, $\Phi(1)=0$ and $\Phi<0$ on $[-1,1)$.

\begin{lemma}\label{lem:series}
There are $p_k\ge0$ with $\sum_kp_k<\infty$ such that $\Phi(c)=\Phi(0)+\sum_{k\ge1}p_kc^k$ on $[-1,1]$, and
$\Phi(C(c))=\chi_1\Phi(c)$ there.
\end{lemma}

\begin{proof}
A composition of power series with nonnegative coefficients has nonnegative coefficients, so
$\Phi_\ell(c)=\Phi_\ell(0)+\sum_{k\ge1}p_k^{(\ell)}c^k$ with $p_k^{(\ell)}\ge0$. Evaluating at $c=1$, the total mass is
$\sum_kp_k^{(\ell)}=-\Phi_\ell(0)\to-\Phi(0)$, so the sequences $p^{(\ell)}$ are bounded in $\ell^1$. Any coefficientwise
limit point $p$ satisfies $\sum_kp_kc^k=\Phi(c)-\Phi(0)$ for $0\le c<1$ by dominated convergence, which determines $p$; so
$p^{(\ell)}\to p$ coefficientwise. By Abel's theorem and the continuity of $\Phi$ at $1$ (the convergence in
Lemma~\ref{lem:cmap} is uniform in $c_0$, as the error there is bounded by its value at $c_0=-1$), $\sum_kp_k=-\Phi(0)$,
which is the limit of the total masses. By Scheff\'e's lemma $p^{(\ell)}\to p$ in $\ell^1$, so $\Phi_\ell$ converges
uniformly on $[-1,1]$ to $\Phi(0)+\sum_kp_kc^k$, which is therefore $\Phi$. The functional equation follows from
$\Phi_{\ell+1}=\chi_1^{-1}\Phi_\ell\circ C$ in the limit.
\end{proof}

\begin{proposition}\label{prop:pk}
Every $p_k$, $k\ge1$, is strictly positive.
\end{proposition}

\begin{proof}
Both sides of the functional equation are power series with nonnegative coefficients beyond the constant term, converging
absolutely on $[0,1]$, so coefficients may be compared:
\begin{equation}\label{eq:pk}
\chi_1p_k=\sum_{j\ge1}p_j\,[c^k]\,C(c)^j\qquad(k\ge1).
\end{equation}
The function $\Phi$ is not constant, since $\Phi(1)=0>\Phi(0)$. It is not a polynomial either: if $p_d>0$ were the last
nonzero coefficient, the right side of \eqref{eq:pk} would contain $p_d\,C^d$, which has infinitely many nonzero
coefficients and cannot be cancelled, while the left side vanishes for $k>d$. Hence $p_j>0$ for infinitely many $j$.
Given $k\ge1$, choose such a $j>k$. Selecting $k$ factors $a_1c$ and $j-k$ factors $a_0$ in $C(c)^j$ gives
$[c^k]\,C^j\ge\binom jk a_0^{j-k}a_1^k>0$, so $p_k\ge\chi_1^{-1}p_j\binom jka_0^{j-k}a_1^k>0$.
\end{proof}

\begin{theorem}\label{thm:positive}
For distinct inputs $x_1,\dots,x_n$ on the fixed-point sphere, $\lambda_{\min}^\perp(-A)>0$.
\end{theorem}

\begin{proof}
The matrix $K=(c_0(x_i,x_j))$ is positive semidefinite with unit diagonal, so $c_0(x_i,x_j)=\langle z_i,z_j\rangle$ for unit
vectors $z_i$ in some $\mathbb R^m$, and they are distinct because $c_0(x_i,x_j)<1$ for $i\ne j$. By
Lemma~\ref{lem:series}, $-A=\Phi(0)J+\sum_{k\ge1}p_kK^{\circ k}$ with Hadamard powers $K^{\circ k}$. On $\mathbf 1^\perp$
the first term vanishes. The kernel $\sum_kp_k\langle z,w\rangle^k$ has infinitely many positive even and infinitely
many positive odd coefficients, by Proposition~\ref{prop:pk}, so it is strictly positive definite on every sphere
\cite{schoenberg,cms}: the matrix $\sum_kp_kK^{\circ k}$ is positive definite, and its quadratic form is bounded below on
the unit sphere of $\mathbf 1^\perp$.
\end{proof}

\section{Inputs off the fixed-point norm}\label{sec:offnorm}
Now let the inputs have arbitrary positive norms. The diagonal recursion $q_{\ell+1}=F(q_\ell)$,
$F(q)=\sigma_w^2\E[\phi(\sqrt q\,z)^2]+\sigma_b^2$, converges to $q^*$, and its rate is $\chi_q=F'(q^*)$.

\begin{lemma}\label{lem:chiq}
$0\le\chi_q<\chi_1$. Consequently $|q_\ell(x)-q^*|=O(\chi_q^\ell)$ for every input $x$.
\end{lemma}

\begin{proof}
By the heat equation, $\frac{d}{dq}\E[f(\sqrt q\,z)]=\frac12\E[f''(\sqrt q\,z)]$. With $f=\phi^2$ and $u=\sqrt{q^*}z$,
\[
\chi_q=\sigma_w^2\,\E\big[\phi'(u)^2+\phi(u)\phi''(u)\big]=\chi_1-2\sigma_w^2\,\E\big[\tanh^2(u)\big(1-\tanh^2(u)\big)\big],
\]
using $\tanh''=-2\tanh\,(1-\tanh^2)$. The last expectation is positive because $q^*>0$, so $\chi_q<\chi_1$; and
$\chi_q\ge0$ because $F$ is increasing, $\tanh^2$ being increasing in $|u|$. Since $F$ is smooth,
$q_{\ell+1}-q^*=\chi_q(q_\ell-q^*)(1+O(|q_\ell-q^*|))$, and the product of the factors converges.
\end{proof}

Put $\rho=\chi_q$. For a pair $x,y$ the correlations obey $c_{\ell+1}=C_\ell(c_\ell)$, where $C_\ell$ is the
correlation map built from the variances $q_\ell(x),q_\ell(y)$. By Lemma~\ref{lem:chiq}, $C_\ell$ and its derivative
differ from $C$ and $C'$ by $O(\rho^\ell)$ uniformly on $[-1,1]$; in particular $C_\ell(1)=1+O(\rho^\ell)$ and
$C_\ell'(1)=\chi_1+O(\rho^\ell)$. The correlations still tend to $1$, because $C$ has $1$ as its only fixed point in
$[-1,1]$ and attracts every point of it.

\begin{lemma}\label{lem:offcmap}
For each pair the limit $A_{xy}=\lim_\ell\chi_1^{-\ell}(1-c_\ell)$ exists and is finite, and
$D_\ell=1-c_\ell=A_{xy}\chi_1^\ell+O\big(\chi_1^{2\ell}+\rho^\ell\big)$.
\end{lemma}

\begin{proof}
Expanding $C_\ell$ at $1$,
\[
D_{\ell+1}=\chi_1D_\ell\big(1+O(D_\ell)\big)+O(\rho^\ell).
\]
Since $D_\ell\to0$, for any $\lambda\in(\chi_1,1)$ the factor $\chi_1(1+O(D_\ell))$ is eventually below $\lambda$, so
$D_\ell=O(\lambda^\ell)$. The factors $1+O(D_\ell)$ then have a convergent product, and comparing with the linear
recursion gives $D_\ell=O\big(\chi_1^\ell+\sum_{k<\ell}\chi_1^{\ell-k}\rho^k\big)=O(\chi_1^\ell)$, because $\rho<\chi_1$.
Dividing by $\chi_1^{\ell+1}$, which is positive since $\phi'>0$,
\[
\chi_1^{-\ell-1}D_{\ell+1}=\chi_1^{-\ell}D_\ell\big(1+O(\chi_1^\ell)\big)+O\big((\rho/\chi_1)^\ell\big),
\]
and both corrections are summable. So $\chi_1^{-\ell}D_\ell$ converges, at the rate of the tails of the two sums.
\end{proof}

The proof of Proposition~\ref{prop:two} carries over with two changes. The diagonal entries $T_\ell(x)=\Theta_\ell(x,x)$
obey $T_{\ell+1}(x)=q_{\ell+1}(x)+\chi_{1,\ell}(x)T_\ell(x)$ with $\chi_{1,\ell}(x)=\chi_1+O(\rho^\ell)$, so for two inputs
$T_\ell(x)-T_\ell(y)$ satisfies a recursion with rate $\chi_1$ and forcing $O(\rho^\ell)$, whence
$T_\ell(x)-T_\ell(y)=O(\chi_1^\ell)$. The gap $E_\ell=T_\ell(x)-\Theta_\ell(x,y)$ satisfies \eqref{eq:gap} up to an
additional forcing $O(\rho^\ell)$, coming from $q_{\ell+1}(x)-\sqrt{q_{\ell+1}(x)q_{\ell+1}(y)}$, from
$\chi_{1,\ell}(x)-\chi_1$ and from $\chi_\ell(c_\ell)-\chi(c_\ell)$. Its contribution to $E_L$ is
$\sum_{\ell<L}\chi_1^{L-1-\ell}O(\rho^\ell)=O(\chi_1^L)$. Hence
\[
E_L=A_{xy}\Big(q^*+\frac{BT^*}{\chi_1}\Big)L\chi_1^L+O(\chi_1^L),
\]
and the Gram matrix satisfies \eqref{eq:gram} with $A=(A_{x_ix_j})$, the differences between the diagonal entries going
into $R_L$. The proof of Proposition~\ref{prop:ninput} uses nothing else, so the asymptotic formula of
Theorem~\ref{thm:main} holds whenever $\lambda_{\min}^\perp(-A)>0$.

The weak inequality survives: the matrix $(c_\ell(x_i,x_j))$ is still positive semidefinite with unit diagonal, and the
argument at the end of Section~\ref{sec:ninput} gives $\lambda_{\min}^\perp(-A)\ge0$. The strict inequality does not
follow from Section~\ref{sec:positive}. There $-A_{ij}=\Phi(c_0(x_i,x_j))$ for one function $\Phi$ given by a series with
positive coefficients, which is what makes the kernel strictly positive definite. Off the sphere each pair is carried by
its own sequence of maps $C_\ell$, and $A$ is no longer a fixed function of the initial correlations.

\section{Numerical checks}\label{sec:numerics}
All Gaussian expectations below are computed by Gauss--Hermite quadrature, the recursions of Section~1 are iterated
exactly in the infinite-width limit, and $\sigma_b^2=0.05$ unless stated otherwise. The scripts are in the
accompanying repository.

\emph{Two inputs.} For two inputs on the fixed-point sphere with initial correlation $c_0$, Proposition~\ref{prop:two}
predicts $E_L/(L\chi_1^L)\to\kappa A$ with a relative error of order $1/L$. Table~\ref{tab:two} lists the computed
ratio $r_L=E_L/(L\chi_1^L\kappa A)$. In every row $(1-r_L)L$ is constant to two digits, as an error $c/L$ requires, and
$\chi_1^{-L}D_L$ agrees with $A$ to five digits by the last depth shown, in line with the geometric rate of
Lemma~\ref{lem:cmap}.

\begin{table}[ht]
\centering
\begin{tabular}{cccccccc}
\hline
$\sigma_w^2$ & $c_0$ & $\chi_1$ & $\kappa$ & $A$ & $L$ & $r_L$ & $(1-r_L)L$\\
\hline
1.5 & 0.0 & 0.93864 & 1.8907 & 0.40464 & 160 & 0.8752 & 20.0\\
    &     &         &        &         & 200 & 0.9001 & 20.0\\
1.5 & 0.5 & 0.93864 & 1.8907 & 0.27864 & 160 & 0.8901 & 17.6\\
    &     &         &        &         & 200 & 0.9121 & 17.6\\
1.2 & 0.0 & 0.84517 & 0.4757 & 0.76500 & 80  & 0.9535 & 3.7\\
    &     &         &        &         & 120 & 0.9690 & 3.7\\
1.2 & 0.5 & 0.84517 & 0.4757 & 0.42715 & 80  & 0.9632 & 2.9\\
    &     &         &        &         & 120 & 0.9755 & 2.9\\
\hline
\end{tabular}
\caption{The two-input gap against the prediction of Proposition~\ref{prop:two}.}\label{tab:two}
\end{table}

\emph{Eight inputs.} Eight random points of the unit sphere in $\mathbb R^5$, rescaled to the fixed-point norm, give a
$8\times8$ matrix $A$ whose restriction to $\mathbf 1^\perp$ is computed from the correlation recursion. The ratio of
$\lambda_{\min}(\Theta_L)/(L\chi_1^L)$ to the prediction $\kappa\lambda_{\min}^\perp(-A)$ of Theorem~\ref{thm:main} is
\begin{gather*}
0.8171,\ 0.9071,\ 0.9443\quad(\sigma_w^2=1.5,\ L=60,120,200),\\
1.0604,\ 1.0311,\ 1.0208\quad(\sigma_w^2=1.2,\ L=40,80,120).
\end{gather*}
Here $|1-\text{ratio}|\,L$ equals $11.0,\ 11.1,\ 11.1$ and $2.4,\ 2.5,\ 2.5$, again an error of order $1/L$. The
error changes sign between the two settings, so the approach to the limit is not monotone in general.

\emph{The series of $\Phi$.} The coefficients $p_k$ of Lemma~\ref{lem:series} were computed by iterating the
composition of truncated power series. At $(\sigma_w^2,\sigma_b^2)=(1.5,0.3)$, $(1.2,0.3)$ and $(1.5,0.05)$ every
$p_k$ with $k\le24$ is positive, the smallest being $1.7\cdot10^{-6}$, $1.5\cdot10^{-7}$ and $5.7\cdot10^{-6}$; the
first coefficients are $p_1=0.372,\ 0.548,\ 0.200$. This agrees with Proposition~\ref{prop:pk}, which needs no
numerical input.

\section*{Acknowledgments and funding}

This research was supported by the European Research Council (ERC) under the European Union's Horizon Europe research and innovation programme (grant agreement No.~101041711), by the Simons Foundation as part of the Collaboration on the Mathematical and Scientific Foundations of Deep Learning, by Heights Labs, by Convex Nexus Capital, by the Israel Science Foundation (grant number 2258/19), and by the Israel Science Foundation (ISF Grant 4101/25).

\section*{Declaration of competing interest}

The author holds equity in Heights Labs and in Convex Nexus Capital and has a management role at Convex Nexus Capital. These entities are included in the funding declaration above.

\section*{Declaration of generative AI and AI-assisted technologies}

Anthropic Claude AI systems were used to develop the mathematical derivations, write and run the numerical checks, perform internal consistency and literature reviews, and write and prepare this manuscript and its publication files. The author selected the problem and directed the work; the author did not carry out the mathematics or the writing. Responsibility for submission or publication rests with the author. The AI-assisted internal reviews, compilation checks and rendered-page inspections are not external peer review or formal verification.

\section*{Data and code availability}

The paper is theoretical. The three scripts that produce the numerical checks of Section~\ref{sec:numerics} are in the \texttt{code} directory of the repository
\[
\texttt{github.com/yspennstate/ntk-smallest-eigenvalue-ordered-phase}.
\]

\begin{thebibliography}{9}
\bibitem{cms} D.~Chen, V.~A.~Menegatto and X.~Sun, \emph{A necessary and sufficient condition for strictly positive
definite functions on spheres}, Proc. Amer. Math. Soc. \textbf{131} (2003), 2733--2740.
\bibitem{xps} L.~Xiao, J.~Pennington and S.~S.~Schoenholz, \emph{Disentangling trainability and generalization in deep
neural networks}, Proc. 37th Int. Conf. Machine Learning (ICML 2020), arXiv:1912.13053.
\bibitem{schoenberg} I.~J.~Schoenberg, \emph{Positive definite functions on spheres}, Duke Math. J. \textbf{9} (1942),
96--108.
\end{thebibliography}


\end{document}
